% ============================================================
\chapter{Introducción}
\label{cap:introduccion}
% ============================================================

La seguridad de una infraestructura Kubernetes queda determinada, en
buena medida, antes de que el clúster llegue a existir: son los manifiestos
declarativos —los ficheros \ac{YAML} que describen el estado deseado de cada
recurso— los que fijan los privilegios de cada carga de trabajo, los
montajes que expone al nodo anfitrión y las relaciones de acceso entre recursos.
Esa configuración se revisa habitualmente recurso a recurso, comprobando cada
manifiesto contra un catálogo de reglas de buenas prácticas. El resultado es una
visión fragmentada del riesgo, porque un compromiso real rara vez depende de una
única configuración incorrecta y sí de la manera en que varias debilidades
menores se encadenan a través de la estructura del clúster.

\section{Motivación}

La adopción masiva de Kubernetes como plataforma de orquestación de
contenedores ha transformado la manera en que las organizaciones despliegan y
gestionan sus aplicaciones. La encuesta anual de 2021 de la \ac{CNCF} recoge que
el 96\,\% de las organizaciones utilizan o evalúan Kubernetes \citep{CNCF2022}.
El informe sectorial \textit{State of Kubernetes Security} de Red Hat de 2024
documenta, por su parte, que el 89\,\% de las organizaciones sufrió al menos un
incidente de seguridad relacionado con Kubernetes durante el último año
\citep{RedHat2024}. Esa proliferación ha ampliado considerablemente la
superficie de ataque: una sola configuración insegura en un manifiesto \ac{YAML} puede habilitar
la escalada de privilegios desde un contenedor comprometido hasta el nodo anfitrión
del clúster o, en casos extremos, hasta la infraestructura subyacente completa
\citep{NSA2022}.

Los enfoques de análisis estático tradicionales —herramientas como Checkov
\citep{Checkov2021} o kube-linter \citep{KubeLinter2021}— detectan
configuraciones incorrectas aisladas en manifiestos individuales. Checkov
incorpora además comprobaciones basadas en grafos sobre las dependencias entre recursos,
pero se trata de reglas declarativas escritas manualmente: ninguna de estas herramientas
dispone de un modelo aprendido de la estructura del clúster ni predice de forma
supervisada la existencia de una cadena de ataque.
Un único indicador de privilegio en un pod —la unidad mínima de despliegue de
Kubernetes, que agrupa uno o varios contenedores— no es suficiente para
determinar si existe una cadena de ataque viable. Lo que habilita rutas de
explotación completas es la combinación de nodos con capacidad de escape, esto
es, de salir del contenedor hacia el nodo anfitrión, con recursos que permiten el
movimiento lateral hacia otras cargas de trabajo del clúster. Esa combinación
sigue la lógica de encadenamiento de etapas del modelo de \textit{kill chain} de
Hutchins, Cloppert y Amin (\citeyear{Hutchins2011}). Las técnicas concretas de
escalada de privilegios en contenedores están catalogadas en el marco MITRE
ATT\&CK for Containers \citep{MITRE2021Containers}.

Herramientas más recientes como KubeHound \citep{KubeHound2023} e
IceKube \citep{IceKube2023} abordan este problema construyendo grafos de
ataque sobre clústeres activos. Requieren, sin embargo, recolectar el estado del
clúster mediante acceso de lectura amplio a su servidor \acs{API} de Kubernetes.
Esa exigencia limita su aplicabilidad como control preventivo en el proceso de
\acs{CICD}, conforme al principio de \textit{shift-left security}, que adelanta
los controles de seguridad a las etapas más tempranas del desarrollo. Este
trabajo aborda ese vacío: construir
un sistema capaz de razonar sobre la \emph{estructura} del clúster a partir de
manifiestos \acs{YAML} estáticos, sin requerir un clúster activo, para identificar
automáticamente cadenas de ataque potenciales antes de que sean desplegadas.

\section{Planteamiento del trabajo}

Formalmente, dado un conjunto de manifiestos \acs{YAML} que describen un clúster
Kubernetes, el problema consiste en:

\begin{enumerate}
  \item Clasificar cada manifiesto individual como seguro o mal configurado (Capa~1).
  \item Modelar el clúster como un grafo heterogéneo y clasificarlo en una de tres
        categorías: \textit{clean} (sin configuraciones erróneas relevantes),
        \textit{isolated} (configuraciones erróneas sin cadena viable) o \textit{attack\_chain}
        (existe una ruta de escalada de privilegios hasta movimiento lateral) (Capa~2).
  \item Combinar las señales de ambas capas, junto con una señal binaria de escape
        derivada de las características de nodo, mediante una función de puntuación
        con pesos optimizados por Algoritmo Genético. El criterio de optimización es
        maximizar la Precisión@5 —la fracción de clústeres de cadena de ataque en el
        top-5 del \textit{ranking}— minimizando los falsos positivos. Sobre el
        \textit{ranking} final se aplica además una restricción de viabilidad
        estructural (Capa~3).
\end{enumerate}

La capacidad de operar sobre manifiestos \acs{YAML} estáticos —sin requerir acceso al
clúster activo— es un requisito de diseño explícito que diferencia el presente trabajo
de las soluciones de grafo de ataque existentes y lo posiciona como un control de
seguridad integrable en \textit{pipelines} \acs{CICD}. En términos generales, el
objetivo es diseñar, implementar y evaluar un sistema de detección de cadenas de
ataque en Kubernetes que complemente los enfoques de análisis estático por
manifiesto individual con el contexto estructural del clúster, mediante la
combinación de aprendizaje automático clásico y aprendizaje sobre grafos. El
sistema opera principalmente sobre manifiestos \acs{YAML} estáticos, si bien
incorpora de forma opcional una vía de consulta directa a un clúster activo. Los
objetivos específicos que
concretan este planteamiento general, junto con la metodología de seis etapas
seguida para alcanzarlos, se detallan en el Capítulo~\ref{cap:objetivos}.

\section{Estructura del trabajo}

El trabajo se organiza en siete capítulos. El Capítulo~\ref{cap:estado} revisa el
estado del arte en seguridad de Kubernetes, herramientas de análisis
estático y grafos de ataque, redes neuronales de grafos y detección de cadenas
de ataque; incluye una tabla comparativa de las herramientas más relevantes y un
análisis del posicionamiento de kubescan. El
Capítulo~\ref{cap:objetivos} define los objetivos específicos y la metodología
de seis etapas. El Capítulo~\ref{cap:requisitos} identifica y justifica los
requisitos funcionales y no funcionales del sistema, delimita su alcance y
describe la estrategia de partición de datos y las métricas de evaluación que
esos requisitos determinan. El Capítulo~\ref{cap:herramienta} describe la
herramienta desarrollada: la arquitectura de tres capas, la construcción del
conjunto de datos y el grafo, el diseño de cada capa, la implementación de kubescan
como herramienta de línea de comandos, la tabla completa de las
28~características del clasificador y la calidad de software y el proceso de
desarrollo seguidos. El Capítulo~\ref{cap:evaluacion} evalúa el
sistema frente a una línea de base del estado del arte, presenta los resultados
de cada capa, sintetiza la métrica de proceso completo y verifica el
cumplimiento de los requisitos, incluida la evaluación de usabilidad con
usuarios expertos. El Capítulo~\ref{cap:conclusiones} recoge las conclusiones
vinculadas a cada objetivo, las limitaciones, las consideraciones éticas y las
líneas de trabajo futuro.
